% Generated draft: requires mathematical and source review.
\documentclass[11pt]{article}
\usepackage[margin=25mm]{geometry}
\usepackage{fontspec}
\setmainfont{Latin Modern Roman}
\usepackage{amsmath,amssymb,mathtools,unicode-math}
\setmathfont{Latin Modern Math}
\usepackage{enumitem,needspace,xurl,hyperref}
\setlength{\parindent}{0pt}
\setlength{\parskip}{5pt}
\setlength{\emergencystretch}{3em}
\newcommand{\N}{\mathbb{N}}
\newcommand{\Z}{\mathbb{Z}}
\newcommand{\Q}{\mathbb{Q}}
\newcommand{\R}{\mathbb{R}}
\newcommand{\C}{\mathbb{C}}
\newcommand{\F}{\mathbb{F}}
\newcommand{\A}{\mathbb{A}}
\newcommand{\PP}{\mathbb{P}}
\newcommand{\E}{\mathbb{E}}
\newcommand{\eps}{\varepsilon}
\newcommand{\dd}{\,\mathrm{d}}
\begin{document}
\section*{UnsolvedMath ID: 30004329}
\textbf{Source problem code:} catalogue entry
\subsection{Definitions and mathematical statement}
Let \(C\subset\PP^2\) be a reduced plane curve, and let \(Z=\operatorname{Sing}(C)\) be its set of singular points, regarded as a reduced set of points. Let \(\pi:X=\operatorname{Bl}_Z(\PP^2)\to\PP^2\) be the blowup at those points, and let \(\widetilde C\) be the strict transform of \(C\), namely the closure of \(\pi^{-1}(C\setminus Z)\) in \(X\). The notation \(\widetilde C^{\,2}\) denotes its self-intersection on \(X\).

The supplied record does not specify the ground field, its characteristic, the geometric-point convention, or the normalization of \(\varepsilon(Z)\). Under the usual interpretation in which \(Z\) consists of distinct geometric points and \(\varepsilon(Z)\) is the unweighted multipoint Seshadri constant of \(\mathcal O_{\PP^2}(1)\), put \(H=\pi^*\mathcal O_{\PP^2}(1)\), and let \(E_z\) be the exceptional curve over \(z\in Z\). Then define
\[
\varepsilon(Z)=\sup\left\{t\in\R_{\geq0}:H-t\sum_{z\in Z}E_z\text{ is nef}\right\}.
\]
Here a real divisor class is nef if its intersection with every integral curve on \(X\) is nonnegative. This makes explicit a conventional interpretation of the notation; the intended field and convention require source verification.

\textbf{Question under this interpretation.} For every reduced plane curve \(C\) as above, does
\[
\widetilde C^{\,2}<-1
\quad\Longrightarrow\quad
\varepsilon(Z)\in\Q
\]
hold? In addition, can \(\varepsilon(Z)\) be computed for such a curve from its singular-point configuration? The source does not specify whether computation means an explicit formula, a terminating exact algorithm, or another form of determination.

The extracted statement concerns arbitrary reduced plane curves, whereas the supplied literature assessment describes a question associated with line arrangements. The intended scope of that assessment relative to this formulation is unresolved.
\subsection{Short English statement}
For a reduced plane curve, blow up all its singular points. If its strict transform has self-intersection less than \(-1\), must the multipoint Seshadri constant of those points be rational, and how can it be computed? The ground field, Seshadri normalization, meaning of computation, and relationship to the supplied line-arrangement assessment require verification.
\subsection{Sources}
\begin{itemize}
\item Supplied local UnsolvedMath registry record, permanent ID \(30004329\), source problem code \texttt{catalogue entry}, title ``Rationality of Seshadri Constants from Plane-Curve Singularities,'' in the Oberwolfach Reports Open Problems collection. No individual proposer is supplied; the proposed year is recorded as 2019.
\item Supplied source citation: \emph{Mini-Workshop: Seshadri Constants} (2020), Oberwolfach Reports, DOI \texttt{10.4171/owr/2019/53}, pp.\ 3271--3313. \url{https://doi.org/10.4171/owr/2019/53}. The supplied background attributes the problem to workshop open-problem discussions. No problem-specific page or passage is supplied.
\item Supplied literature source: ``OWR report on Seshadri constants.'' \url{https://ems.press/content/serial-article-files/46833}.
\item Literature assessment dated exactly \texttt{2026-08-22}: the supplied curation labels the problem open and states that general computation and rationality questions for multipoint Seshadri constants associated with line arrangements are not settled by the known examples and bounds. This is supplied dated curation, not an independently verified literature assessment. The background expressly requests independent verification of time-sensitive status and cited claims.
\item The registry labels the statement \texttt{corrected\_verified}, with an accompanying explanation that only formula delimiters, notation, grammar, and typography were repaired and that mathematical content and scope were preserved. This label is supplied metadata and does not establish mathematical correctness or independent source verification.
\end{itemize}
\end{document}
